
Or
@@M0@@

Par
ailleurs, @@M1@@~; nous noterons @@M2@@ la classe dans @@M3@@ de @@M4@@. La relation (1.11.1) donne donc finalement :
@@M5@@

On considère alors le complexe :
@@M6@@
dont les flèches sont définies par :
@@M7@@
Si on considère @@M8@@ comme @@M9@@-Module gradué réduit au degré @@M10@@, c'est un complexe de @@M11@@-Modules gradués, les morphismes étant eux-mêmes gradués.

Ce complexe est une résolution de @@M12@@ : sur un ouvert où @@M13@@ est libre, @@M14@@ s'identifie à une Algèbre de polynomes @@M15@@, et le complexe au complexe de Koszul défini par la suite @@M16@@-régulière @@M17@@; on sait que ce dernier est une résolution de @@M18@@.

En écrivant que chacune des composantes homogènes du complexe de degré @@M19@@ est acyclique, on obtient donc des relations dans @@M20@@ :
@@M21@@
(en posant @@M22@@ pour @@M23@@).

Supposons la relation (1.11.2) vraie pour @@M24@@; nous allons montrer qu'elle est encore vraie pour @@M25@@, ce qui permettra de redescendre par récurrence jusqu'à @@M26@@ (nous verrons plus loin que cela suffit à entrainer la nullité des @@M27@@).

